\chapter{Cross-Modal Evidence}
\label{ch:essay-crossmodal}

The formal apparatus developed in the preceding chapter applies whenever structure is projected through a physical channel. Language, gesture, and music differ in which aspects of intentional structure each modality preserves and which it discards, but the projection problem remains the same across all three: determining what survives the transformation from latent structure to observable signal. A worked example illustrates the application of Representation Invariance before turning to measured evidence from the repository's implemented experiments.

Consider the communicative act of requesting transfer of an object, realized in two modalities. In speech, this might be expressed as the utterance ``give me the book,'' consisting of a verb, two pronouns, and a definite noun phrase linearized according to the grammatical constraints of English. In gesture, the same intention might be realized as a sequence consisting of a point toward the speaker, a grasping motion, and a point toward the object, with no single element corresponding uniquely to the verb or to definiteness. The question is whether these two realizations count as representationally invariant in the sense established by Theorem~1.1.

Establishing invariance rigorously requires constructing an explicit admissibility-isomorphism $\varphi$ between the two state spaces and verifying that admissible continuations in one modality correspond exactly to admissible continuations in the other. The speech realization admits continuations such as ``give it to me'' or ``hand me the book,'' which preserve the request structure while varying lexical or syntactic detail. The gesture realization admits continuations such as repeating the sequence with greater emphasis or adding an upward palm orientation to mark urgency. The question is whether there exists a map $\varphi$ such that every admissible speech continuation maps to an admissible gesture continuation and conversely, in both directions, preserving the full reachable-set structure. Without exhibiting such a map and verifying the biconditional, the claim that the two realizations are invariant remains illustrative rather than proven. What can be stated without constructing $\varphi$ is that if such an isomorphism exists, then Corollary~1.2 applies: substrate-specific properties such as the acoustic-versus-motor distinction, production latency, or articulatory cost are not semantic properties unless they constrain the admissibility relation itself.

The question of what information survives projection and what is lost has been investigated directly in the repository's implemented experiments, though not yet across all modalities. The measurements available as of this writing come from experiments on linguistic evolution rather than cross-modal comparison, and they should be understood as evidence concerning projection within a single modality rather than across modalities.

The Phonological Drift experiment, implemented in \texttt{experiments/phonological\_drift.py}, models sound change as a population process rather than as a predetermined historical sequence. Twenty-five speakers occupy a geographical grid, and probabilistic sound changes spread through neighbor influence. Over fifteen generations, the experiment tracks lexical divergence between speakers who began with identical pronunciations. A representative run produces an average lexical divergence of $0.377$ and a maximum divergence of $0.750$ between the most distant speaker pairs. Two speakers pronouncing the ancestral word \textit{pater} may produce \textit{faθer} and \textit{pater} respectively after fifteen generations, despite both trajectories consisting entirely of admissible phonological changes at each step. The measured divergence quantifies how much phonological structure has been altered by the projection through time and population structure, though the experiment does not yet measure how much of the ancestral state remains reconstructible from the divergent outcomes.

The Lexical Natural Selection experiment, implemented in \texttt{experiments/lexical\_selection.py}, models competition among synonymous forms without reducing fitness to a universal scalar. Variants differ along dimensions including length, articulatory cost, regularity, and prestige, with relative advantages depending on context. Four meanings begin with three competing variants each and evolve over one hundred fifty time steps. A representative run for the meaning \textit{father} produces a distribution in which the variant \textit{pa} accounts for $36.5\%$ of usage, \textit{papa} for $34.8\%$, and \textit{pater} for $28.7\%$. The experiment demonstrates that selection can alter the frequency distribution of competing forms without eliminating variation entirely, though again it does not yet measure how much of the selection history remains recoverable from the final distribution alone.

The Reconstruct experiment, implemented in \texttt{experiments/reconstruct.py}, closes the experimental loop by attempting historical inference without access to the simulator's privileged ground truth. A proto-language is evolved into four daughter languages over twenty time steps, the ancestral state is concealed, and a simplified comparative method attempts reconstruction using only the surviving daughters. A representative run achieves correct reconstruction for six of ten test items, partial reconstruction for four items, and complete failure for zero items, yielding a sixty-percent exact-match accuracy. This is the first direct measurement of reconstruction accuracy under the complete $\History \to \Observable \to \Reconstruction$ protocol, and it demonstrates that some information survives the projection from history to observable well enough to permit reconstruction while other information is lost.

The Recoverability experiment, implemented in \texttt{experiments/recoverability.py}, investigates the distinction between reconstruction error and fundamental non-identifiability. Rather than asking only whether a particular history can be reconstructed, the experiment generates multiple independent histories and searches for cases in which $H_1 \neq H_2$ but $\obsmap(H_1) \approx \obsmap(H_2)$. A representative run across ten generated histories identifies twenty-four pairs of distinct histories with similar observable outcomes, including cases where observable distance is measured as zero despite the histories differing in event count and evolutionary trajectory. These cases instantiate the structural direction of Theorem~4.2: distinct histories sharing the same reachable-set structure produce indistinguishable observables, and no reconstruction procedure operating solely on observables can recover which history actually occurred.

The measurements reported above are all drawn from experiments operating within a single modality, specifically linguistic evolution under sound change, lexical competition, and historical divergence. Cross-modal measurements comparing information loss across language, gesture, music, and audio are not yet available from implemented experiments. Claims about the percentage of trajectory information lost when continuous gesture is sampled to discrete frames, or the percentage of semantic content lost when structured meaning is linearized into speech, remain hypotheses to be tested rather than measured results. A proper cross-modal comparison would require implementing projectors and inverse engines for each modality as specified in the unified framework, running controlled experiments in which the same semantic structure is projected through different channels, and measuring what fraction of the original structure remains recoverable in each case. Such experiments would test whether the information-loss patterns predicted by the four theorems manifest differently across modalities or whether the projection problem exhibits modality-independent structure.

The evidence available from the repository's current experiments demonstrates that the $\History \to \Observable \to \Reconstruction$ framework can distinguish three questions that are easily conflated: what evolutionary history actually occurred, what evidence of that history survived to the present, and what an observer can legitimately infer from that surviving evidence. The Reconstruct experiment shows that reconstruction can succeed partially even when information has been lost. The Recoverability experiment shows that distinct histories can become observationally indistinguishable, instantiating the collision condition required by Theorem~4.2. Together, these results establish that the framework's epistemic distinctions are not merely definitional but measurable, and that the theorems proven in \emph{Admissible Reconstruction} apply to the implemented experiments. What remains to be measured is whether the same distinctions apply with comparable quantitative characteristics when the projection operates across modalities rather than solely through time within a single modality.
